% TOP_PROBLEM: {"problemId": "problem.hall-s-conjecture-on-gaps-between-squares-and-cubes", "canonicalTitle": "Hall's Conjecture on Gaps Between Squares and Cubes", "releaseRank": 465, "primaryDomain": "number_theory_arithmetic_geometry", "primaryDomainLabel": "Number theory & arithmetic geometry", "displayStatus": "open", "releaseStatus": "open"}
% !TeX program = lualatex
% ATTEMPT_STATUS: unresolved
\documentclass[11pt]{article}
\usepackage[margin=25mm]{geometry}
\usepackage{fontspec}
\setmainfont{Latin Modern Roman}
\usepackage{amsmath,amssymb,mathtools}
\usepackage{unicode-math}
\setmathfont{Latin Modern Math}
\usepackage{xurl,hyperref}
\hypersetup{colorlinks=true,urlcolor=blue}
\setcounter{secnumdepth}{0}
\setlength{\parindent}{0pt}
\setlength{\parskip}{5pt}
\setlength{\emergencystretch}{3em}
\begin{document}
\section*{\texorpdfstring{Hall's Conjecture on Gaps Between Squares and Cubes}{Hall's Conjecture on Gaps Between Squares and Cubes}}\label{465}
\subsection{Definitions and mathematical statement}
For positive integers $x,y$, exclude exact equality $y^2=x^3$. Hall's selected small-power-loss assertion is
\[
 \forall{\varepsilon}>0\ \exists c_{\varepsilon}>0\ \forall x,y\in{\mathbb{N}}:\quad
 y^2\ne x^3\Longrightarrow
 |y^2-x^3|>c_{\varepsilon} x^{1/2-{\varepsilon}}.
\]
The constant depends on ${\varepsilon}$ only, not on the particular square and cube. The absolute value allows either one to be larger. This formulation retains the arbitrary ${\varepsilon}$-loss; the claim with exponent exactly $1/2$ and one positive uniform constant is stronger and is not the displayed target. A finite collection of unusually close pairs cannot disprove a statement permitting an unspecified positive constant.
\par\smallskip\textit{Formulation and concept references:} \href{https://doi.org/10.1007/BF01140190}{[S1]}.

\subsection{Short English statement}
Apart from equality, must an integer square stay at least a constant times a nearly square-root power of x away from the cube of x?
\subsection{Research attempt}
\paragraph{Scope and literature check (27 September 2026).}
The selected statement includes an arbitrary positive loss in the exponent, permits both signs of $y^2-x^3$, and does not impose coprimality. These details matter. The historical Hall reference [S1] is retained as supplied; its full text was not available in this check. Brian Conrad's primary lecture [E1] discusses the distinction between exponent variants and the conditional connection with the $abc$ conjecture. The attempt below proves that implication explicitly without a coprimality restriction, proves a polynomial analogue from a derivative argument, and checks exact close square--cube pairs. It supplies no unconditional integer lower bound of the requested strength.

\paragraph{Elementary reductions and the actual difficult range.}
Put $d=|y^2-x^3|>0$. Integrality gives $d\geq1$. For $\varepsilon\geq1/2$, the requested estimate already follows with $c_\varepsilon=1/2$, since $x^{1/2-\varepsilon}\leq1$. Thus only $0<\varepsilon<1/2$ requires new information.

If $x=a^2$ is a square, then $x^3=a^6$ is itself a square. Among positive integers $y\ne a^3$, the distance from $y^2$ to $a^6$ is at least $2a^3-1$; for $a=1$ the positive-neighbor restriction only increases this lower bound. Hence
\[
                        |y^2-x^3|\geq2x^{3/2}-1
\]
in that subcase, far stronger than required. The challenge is the approach of a nonsquare cube to a square.

For a fixed nonsquare $x$, the smallest absolute difference is attained at $y=\lfloor x^{3/2}\rfloor$ or $y=\lceil x^{3/2}\rceil$: outside these two choices the square is farther from $x^3$ by monotonicity. The identity
\begin{equation}\label{a465:distance}
                 |y-x\sqrt{x}|=\frac{|y^2-x^3|}{y+x\sqrt{x}}
\end{equation}
shows that in the near-cancellation range a lower bound of order $x^{1/2-\varepsilon}$ for $d$ requires a lower bound of order $x^{-1-\varepsilon}$ for the distance of $x\sqrt x$ from an integer. Merely saying $\sqrt x$ is irrational gives a nonzero distance, with no uniform lower bound at this scale. A Diophantine approximation theorem for one fixed quadratic irrational also cannot be used with a constant independent of the varying $x$ without proving that uniformity.

\paragraph{Why scaling a good pair cannot disprove the conjecture.}
Under the natural transformation $(x,y)\mapsto(t^2x,t^3y)$,
\[
 |(t^3y)^2-(t^2x)^3|=t^6d,
 \qquad
 \frac{t^6d}{(t^2x)^{1/2-\varepsilon}}
       =t^{5+2\varepsilon}\frac d{x^{1/2-\varepsilon}}.
\]
Thus dilating a close pair makes the normalized ratio larger. This is one reason a single spectacular numerical example supplies no contradiction. To disprove the selected statement one needs a fixed $\varepsilon>0$ and a sequence with $x\to\infty$ and $d/x^{1/2-\varepsilon}\to0$. A finite collection only bounds how large an admissible constant might be; it cannot exclude all positive constants.

Nor can one simply divide $x$ and $y$ by their ordinary greatest common divisor to reach the coprime case: the powers two and three scale differently. The conditional argument below removes the common divisor at the level of the three terms in an additive equation, where the cancellation is exact.

\paragraph{Proposition 465.1: the standard $abc$ conjecture implies this statement.}
Assume the following conjectural input: for each $\delta>0$ there is $C_\delta>0$ such that coprime positive integers $A+B=C$ satisfy
\begin{equation}\label{a465:abc}
                C\leq C_\delta\operatorname{rad}(ABC)^{1+\delta}.
\end{equation}
Here $\operatorname{rad}(m)$ is the product of the different primes dividing $m$, with $\operatorname{rad}(1)=1$. This hypothesis is not treated as an available unconditional theorem.

Let $g=\gcd(x^3,y^2)$ and $M=\max\{x^3,y^2\}$. Since $g$ divides $d$, the three integers
\[
 A=\frac{\min\{x^3,y^2\}}g,\qquad B=\frac dg,
 \qquad C=\frac Mg
\]
are positive and pairwise coprime, and satisfy $A+B=C$. A prime dividing $A$ or $C$ divides $xy$. Consequently
\[
 \operatorname{rad}(ABC)
 \leq\operatorname{rad}(xy)\operatorname{rad}(d/g)
 \leq xy\frac dg.
\]
Applying \eqref{a465:abc} and multiplying by $g$ gives
\begin{equation}\label{a465:gcd}
 M\leq C_\delta(xyd)^{1+\delta}g^{-\delta}
       \leq C_\delta(xyd)^{1+\delta}.
\end{equation}
The common divisor improves this estimate, so no coprimality assumption on $x,y$ is needed.

If $x^3/2\leq y^2\leq2x^3$, then $xy\leq\sqrt2x^{5/2}$ and $M\geq x^3$. Rearranging \eqref{a465:gcd} yields
\begin{equation}\label{a465:conditional}
 d\geq 2^{-1/2}C_\delta^{-1/(1+\delta)}
       x^{3/(1+\delta)-5/2}
   =c'_{\delta}x^{1/2-3\delta/(1+\delta)}.
\end{equation}
Take $\delta=\varepsilon/3$. Then $3\delta/(1+\delta)<\varepsilon$, so for $x\geq1$ the right side is at least $c'_{\delta}x^{1/2-\varepsilon}$. Outside the comparable range, $d>x^3/2$ or $d>x^3$, which is stronger still. Choosing a positive constant smaller than both $c'_{\delta}$ and $1/2$ turns the non-strict inequalities into the strictly greater inequality in the problem, uniformly over both cases.

This completes the conditional proof, including the signs, the gcd, and the dependence of constants. Its only unproved ingredient is explicitly \eqref{a465:abc}; discarding that label would falsely present the conditional argument as a solution.

\paragraph{A polynomial analogue that can be proved unconditionally.}
Over $\mathbb C[t]$, differentiation supplies an additional tool. Suppose nonzero pairwise coprime polynomials $A,B,C$ satisfy $A+B=C$ and are not all constant. Let $R$ be the squarefree product of all roots occurring in $ABC$, so $\deg R$ counts the distinct roots. Then
\begin{equation}\label{a465:mason}
              \max\{\deg A,\deg B,\deg C\}\leq\deg R-1.
\end{equation}
Here is the complete degree argument. The Wronskian $W=A'B-AB'$ is nonzero: otherwise the derivative of $A/B$ vanishes, hence $A/B$ is constant in characteristic zero, and coprimality forces both $A,B$, and then $C$, to be constant. At a root of $A$ of multiplicity $r$, $W$ vanishes to order at least $r-1$; the same holds for $B$. Since
\[
                      W=A'C-AC',
\]
the corresponding statement holds for roots of $C$. Pairwise coprimality makes those root sets disjoint. Thus $(ABC)/R$ divides $W$. Comparing degrees gives
\[
 \deg A+\deg B+\deg C-\deg R
       \leq\deg W\leq\deg A+\deg B-1,
\]
so $\deg C\leq\deg R-1$. Permuting the equation, with signs as needed, proves the other two bounds. This proves \eqref{a465:mason}, the polynomial $abc$ or Mason--Stothers inequality, in the setting needed here.

Apply it to coprime nonconstant polynomials $X,Y$ with $\deg X=2m$, $\deg Y=3m$, $m\geq1$, and $D=X^3-Y^2\ne0$. The polynomials $X^3,-Y^2,D$ are pairwise coprime, since a root shared with $D$ would be shared by $X$ and $Y$. Their distinct roots number at most $2m+3m+\deg D$, while their maximum degree is at least $6m$. Hence
\begin{equation}\label{a465:polygap}
                              \deg D\geq m+1.
\end{equation}
This rules out polynomial families with these degree and coprimality hypotheses whose cancellation leaves degree at most $m$. It is an unconditional structural obstruction, not an integer lower bound.

For example,
\[
 X=t^2+2,\qquad Y=t^3+3t,
 \qquad X^3-Y^2=3t^2+8.
\]
The two polynomials are coprime: a root of $X$ has $t^2=-2$, hence $Y=t\ne0$. The degrees are two and three, and the residual degree two attains \eqref{a465:polygap} for $m=1$. On large positive integer substitution, this residual is comparable to $X$, rather than to $X^{1/2}$. It therefore is not a near-counterexample at the conjectured scale.

The derivative proof cannot be transferred just by replacing degree with logarithmic size. Integers have no derivative satisfying the multiplicity reduction used in $W$, and the crucial nonzero auxiliary polynomial and its degree bound would need an arithmetic substitute. The missing substitute is the substance of the analogy, not a routine final step. Also, a polynomial family is a special form of an infinite sequence of integer pairs; ruling out certain families would not rule out all possible pairs.

\paragraph{What fixed-residual finiteness does and does not say.}
Finiteness of the integer solutions to each fixed equation $y^2=x^3+k$, $k\ne0$, is compatible with solutions whose sizes grow extremely rapidly as $|k|$ grows. Hall's assertion is a uniform quantitative relationship as the residual varies. A separate finiteness statement for every $k$ supplies no bound with a common exponent or constant in $k$.

More explicitly, for $0<\varepsilon<1/2$ the proposed lower bound would imply
\[
                 x<\left(\frac{|k|}{c_\varepsilon}\right)^{1/(1/2-\varepsilon)}
                    \quad\text{when }y^2=x^3+k.
\]
As $\varepsilon\downarrow0$ the exponent approaches two from above. An estimate exponential in a positive power of $|k|$ can prove effective finiteness and still be far too weak to imply any one of these polynomial estimates. This explains why a general algorithm to solve a fixed Mordell equation is not itself a proof of the selected uniform bound.

\paragraph{Exact numerical counterchecks.}
The diagnostic enumerates $1\leq x\leq100000$. It computes $r=\lfloor\sqrt{x^3}\rfloor$ by integer square root and tests the positive candidates among $r-1,r,r+1$, discarding zero residual. This contains the nearest admissible square in every case, including square $x$, where the excluded exact square must be replaced by a neighbor. The scan records new minima of $d^2/x$, a rational quantity that compares $d/\sqrt x$ without rounding. It keeps the sign of $y^2-x^3$ separately.

Two further exact checks illustrate scale without making an asymptotic claim:
\[
 378661^2-5234^3=17,
\]
\[
 447884928428402042307918^2
   -5853886516781223^3=-1641843.
\]
The second pair is the Elkies example recorded in [E1]; the displayed equality is independently checked by integer arithmetic. The scan is not described as a record search over all known pairs. It also verifies the polynomial identity above by coefficient arithmetic and checks the scaling identity on a grid of inputs.

A small value of $d/\sqrt x$ constrains the strong exponent-$1/2$ constant, if such a constant exists. For the selected weaker statement the ratio is $d x^{\varepsilon}/\sqrt x$, so fixing a positive $\varepsilon$ changes what would constitute a counterexample sequence. Neither displayed pair, nor the finite scan, supplies that sequence.

\paragraph{Unresolved step.}
The unconditional results here cover square $x$, the trivial exponent range, exact reductions, scaling obstructions, and a fully proved polynomial analogue. The $abc$ deduction would establish the complete stated bound with all its quantifiers, but it remains conditional on \eqref{a465:abc}. The integer problem requires a uniform lower bound for the near-cancellation in \eqref{a465:distance} beyond integrality. No such unconditional estimate is proved in this attempt, and the selected Hall conjecture remains unresolved here.

\textbf{UNRESOLVED.} The partial results above do not establish the full selected statement.

\subsection{Sources}
\begin{itemize}
\item[S1]M. Hall, The Diophantine equation x\textasciicircum{}3-y\textasciicircum{}2=k (1971).
\url{https://doi.org/10.1007/BF01140190}.
\textit{Formulation source.}
\item[E1] B. Conrad, ``The ABC conjecture,'' Stony Brook colloquium lecture, 12 September 2013.
\url{https://www.math.stonybrook.edu/Videos/Colloquium/PDFs/20130912-Conrad.pdf}.
\textit{Primary exposition of Hall variants, the Elkies numerical example, and the conditional $abc$ implication; the argument here is rederived and remains conditional. Checked 27 September 2026.}
\end{itemize}
\end{document}
